\scanpage{31}
% @ 1
\uline{Proposition} $G$ \deleted{\unsure{algebraic group}} \tikz[baseline=(s.base)]{\node[inner sep=1pt](s){affine};\node[anchor=south east,inner sep=1pt](a) at ([xshift=2pt,yshift=7pt]s.north west){\inserted{group \inserted{(monoid)}}};\draw ([xshift=-8pt,yshift=-4pt]s.south west) to[out=90,in=190] ([xshift=1pt,yshift=-2pt]a.south west)--([xshift=1pt,yshift=-2pt]a.south east) to[out=0,in=-90] ([xshift=3pt,yshift=3pt]a.south east);} over the field $k$. Then

% @ 2–3
a) $G$ \deleted{\unsure{G}} is a projective limit \tikz[baseline=(s.base)]{\node[inner sep=1pt](s){of};\node[anchor=south east,inner sep=1pt](a) at ([xshift=2pt,yshift=7pt]s.north west){\inserted{\unsure{strict filtering|filtering of|filtering}}};\draw ([xshift=-8pt,yshift=-4pt]s.south west) to[out=90,in=190] ([xshift=1pt,yshift=-2pt]a.south west)--([xshift=1pt,yshift=-2pt]a.south east) to[out=0,in=-90] ([xshift=3pt,yshift=3pt]a.south east);} affine \tikz[baseline=(s.base)]{\node[inner sep=1pt](s){(monoids)};\node[anchor=south east,inner sep=1pt](a) at ([xshift=2pt,yshift=7pt]s.north west){\inserted{groups}};\draw ([xshift=-8pt,yshift=-4pt]s.south west) to[out=90,in=190] ([xshift=1pt,yshift=-2pt]a.south west)--([xshift=1pt,yshift=-2pt]a.south east) to[out=0,in=-90] ([xshift=3pt,yshift=3pt]a.south east);} of finite type over $k$, \deleted{\unsure{and}} \deleted{\unsure{b} with morphisms}

% @ 4–5
b) If $G$ is an \tikz[baseline=(s.base)]{\node[inner sep=1pt](s){affine};\node[anchor=south east,inner sep=1pt](a) at ([xshift=2pt,yshift=7pt]s.north west){\inserted{(\unsure{group|group})}};\draw ([xshift=-8pt,yshift=-4pt]s.south west) to[out=90,in=190] ([xshift=1pt,yshift=-2pt]a.south west)--([xshift=1pt,yshift=-2pt]a.south east) to[out=0,in=-90] ([xshift=3pt,yshift=3pt]a.south east);} monoid of finite type over $k$, one can find an \deleted{\unsure{i|immersion}} integer $n$ and a closed immersion
% @ 6
\[
 G\mathrel{\tikz[baseline=-.3ex,x=1em,y=1em]{\draw[line width=.45pt] (0,0)--(1,0) (1,-.12)--(1,.12);}}\tikz[baseline=(s.base)]{\node[inner sep=1pt](s){\(\deleted{\unsure{\mathrm{GL}}}M(n)_{k}^{\unsure{\times|x}}\)};\node[anchor=north,inner sep=1pt](a) at ([xshift=12pt,yshift=-15pt]s.south){\inserted{scheme of $n\times n$ matrices}};\draw[->] ([xshift=3pt,yshift=-1pt]s.south) to[out=-100,in=170] ([xshift=-12pt,yshift=2pt]a.north);}.\quad
 (\text{resp. }G\to\mathrm{GL}(n))
\]

% @ 7–9
\uline{Proof} $G\longleftrightarrow A$, \deleted{\unsure{and}} and $A=\underset{i}{\underset{\mathrel{\tikz[baseline=-.3ex,x=1em,y=1em]{\draw[line width=.45pt] (0,0)--(1,0) (1,-.12)--(1,.12);}}}{\mathrm{lim}}}\,A_i$, $A_i$ the subalgebras of finite type of $A$. We are interested in the subalgebras $B$ which are ``stable'' under $\unsure{\Pi|\pi}:A\otimes A$, now it is clear that if
% @ 10–12
$M\subset A$ is stable under $\Pi$, then also the subalgebra \tikz[baseline=(s.base)]{\node[inner sep=1pt](s){which it};\node[anchor=south east,inner sep=1pt](a) at ([xshift=2pt,yshift=7pt]s.north west){\inserted{$B$}};\draw ([xshift=-8pt,yshift=-4pt]s.south west) to[out=90,in=190] ([xshift=1pt,yshift=-2pt]a.south west)--([xshift=1pt,yshift=-2pt]a.south east) to[out=0,in=-90] ([xshift=3pt,yshift=3pt]a.south east);} generates. Thus corollary 2 is for us that one can assume the $A_i$ stable under $\Pi$. If $G$ a group, whence an antipode
% @ 13–17
$\sigma$ of $A$, then if $M$ stable under $\sigma$, then the algebra $B$ also. Now if $M$ stable under $\Pi$, $\sigma M$ also [$M$ left submodule \tikz[baseline=(s.base)]{\node[inner sep=1pt](s){over};\node[anchor=south east,inner sep=1pt](a) at ([xshift=2pt,yshift=7pt]s.north west){\inserted{(right)}};\draw ([xshift=-8pt,yshift=-4pt]s.south west) to[out=90,in=190] ([xshift=1pt,yshift=-2pt]a.south west)--([xshift=1pt,yshift=-2pt]a.south east) to[out=0,in=-90] ([xshift=3pt,yshift=3pt]a.south east);} $\check A$ $\Longleftrightarrow$ $\sigma M$ right (left) submodule over $\check A$] thus $M+\sigma M$ also, thus one can assume the $A_i$ moreover stable under $\sigma$, thus the $G_i$ are groups.

% @ 18–21
Suppose \unsure{now|now} $A$ of finite type, thus we have a $k$-submodule \tikz[baseline=(s.base)]{\node[inner sep=1pt](s){$M$};\node[anchor=south east,inner sep=1pt](a) at ([xshift=2pt,yshift=7pt]s.north west){\inserted{of finite type}};\draw ([xshift=-8pt,yshift=-4pt]s.south west) to[out=90,in=190] ([xshift=1pt,yshift=-2pt]a.south west)--([xshift=1pt,yshift=-2pt]a.south east) to[out=0,in=-90] ([xshift=3pt,yshift=3pt]a.south east);} of $A$ which generates $A$ as an algebra. \deleted{Consider} One can assume $M$ \unsure{stable|stable} under $\check A$ on the \tikz[baseline=(s.base)]{\node[inner sep=1pt](s){\deleted{\unsure{right}}};\node[anchor=south,inner sep=1pt] at ([yshift=3pt]s.north){\deleted{\unsure{left|left}}};\node[anchor=north,inner sep=1pt] at ([yshift=-3pt]s.south){\inserted{right}};} thus that $G$ acts on it on the \deleted{\unsure{left}} \inserted{right}. Thus
% @ 22
\[
 G\longrightarrow\deleted{\unsure{\mathrm{GL}|{\mathrm{End}}}}\,\underline{\mathrm{Aut}}_{k}(M)\quad(\simeq\mathrm{GL}(n)_{k})
\]
% @ 23–24
With $n=\dim M$. Let us prove that it is a monomorphism. \deleted{\unsure{We have}} Let \deleted{$\unsure{g}\ \varepsilon(f)\ \unsure{=0}$} Consider \tikz[baseline=(s.base)]{\node[inner sep=1pt](s){the};\node[anchor=south east,inner sep=1pt](a) at ([xshift=2pt,yshift=7pt]s.north west){\inserted{for \unsure{every|$\forall$} $\varphi\in M$}};\draw ([xshift=-8pt,yshift=-4pt]s.south west) to[out=90,in=190] ([xshift=1pt,yshift=-2pt]a.south west)--([xshift=1pt,yshift=-2pt]a.south east) to[out=0,in=-90] ([xshift=3pt,yshift=3pt]a.south east);} function on $G$
% @ 25
\[
 G\longrightarrow\mathbb E\qquad \deleted{\unsure{g}}\quad
 g\longmapsto\varepsilon(g.\varphi)=(g\varphi)(1)=\varphi(g)
\]
% @ 26–28
\deleted{\unsure{thus} $\unsure{E|k}'$} Thus \deleted{\unsure{g}} $M\unsure{\subset|\in}\mathrm{Im}\,[\text{ring }\tikz[baseline=(s.base)]{\node[inner sep=1pt](s){of};\node[anchor=south east,inner sep=1pt](a) at ([xshift=-2pt,yshift=7pt]s.north west){\inserted{affine}};\draw ([xshift=-12pt,yshift=-2pt]s.north west) to[out=100,in=0] ([xshift=-10pt,yshift=-2pt]a.south east);\draw ([xshift=-12pt,yshift=-2pt]s.north west)--([xshift=-2pt,yshift=10pt]s.north west);}\,\underline{\mathrm{End}}_k(M)\mathrel{\tikz[baseline=-.3ex,x=1em,y=1em]{\draw[line width=.45pt] (0,0)--(1,0) (1,-.12)--(1,.12);}}A]$ \deleted{\unsure{thus}} \unsure{The affine ring|the affine ring} of $\underline{\tikz[baseline=(e.base)]{\node[inner sep=0pt](e){\(\mathrm E\)};\draw ([xshift=1pt,yshift=2pt]e.north west) to[out=20,in=140] ([xshift=2pt,yshift=0pt]e.north east);}\mathrm{nd}}_k(M)$ and a fortiori that of $\mathrm{GL}(n)$ goes \uline{onto} $A$, OK.
% archive sheet number: 30
